\chapter{Beurteilung}
\label{cha:Beurteilung}

Die Beurteilung der Arbeit erfolgt nach einem Schema, welches die Hauptkriterien
\begin{itemize}
    \item Arbeitsstil,
    \item Ergebnisse,
    \item Ausarbeitung und
    \item Abschlussvotrag
\end{itemize}
umfasst.
Im Folgenden sind kurz die wesentlichen Aspekte, die darunter jeweils betrachtet werden, beschrieben.


\paragraph{Arbeitsstil}
Hierunter fällt die Bewertung des (hoffentlich) systematischen und methodischen Vorgehens.
Dies umfasst neben der sinnvollen zeitlichen Vorgehensweise auch die Literaturrecherche zu Beginn der Arbeit sowie die Kooperation und Kommunikation mit dem Betreuer.
Daneben spielt auch die Selbständigkeit sowie das fachliche Verständnis eine große Rolle.

Bei Gruppenarbeiten wird auch die Zusammenarbeit der Gruppe bewertet.
Dies meint nicht unbedingt, wie harmonisch eine Gruppe zusammenarbeitet, sondern ob die Aufgabe in sinnvolle Pakete zerlegt wurde, so dass durch die Bearbeitung in der Gruppe ein Mehrwert im Vergleich zu Einzelarbeiten erreicht wurde.


\paragraph{Ergebnisse}
Hierbei wird bewertet, inwieweit die Arbeit vollständig ist und welche Qualität die Ergebnisse haben.
\Dah sind alle geforderten Punkte (und \ggf naheliegende Aspekte) berücksichtigt und in einer solchen Form bearbeitet, dass diese direkt weiterverwendet werden können?


\paragraph{Ausarbeitung}
Ein wesentliches formales Kriterium an dieser Stelle ist die korrekte und vollständige Zitierweise.
Schwerwiegende Fehler an dieser Stelle (Plagiate) führen zu einem Nicht-Bestehen der gesamten Arbeit!

Darüber hinaus sind die wesentlichen Kriterien die Gliederung der Arbeit sowie die Vollständigkeit und Verständlichkeit.
\Dah ist der Aufbau und Argumentationsweg der Arbeit nachvollziehbar, sind die Begründungen schlüssig?

Die sprachliche Ausdrucksweise (über die Verständlichkeit hinaus) und Korrektheit, sowie die prinzipielle Sorgfalt der Gestaltung (im Wesentlichen Plots und Zeichnungen) werden hierbei auch berücksichtigt, aber mit einem geringeren Gewicht.


\paragraph{Abschlussvortag}
Auch hier wird die Gliederung der Folien bewertet.
Wurde sich auf das wichtigste beschränkt und folgt der Vortrag einem roten Faden?
Auch die Gestaltung der Folien spielt hier eine Rolle bei der Bewertung.
Der Vortragsstil wird ebenfalls beachtet, jedoch mit einem etwas geringeren Gewicht.

Neben dem eigentlichen Vortrag spielt hier auch die Diskussion eine wesentliche Rolle bei der Bewertung.
\Dah wurden die Fragen korrekt erfasst und entsprechend eine überzeugende Antwort darauf gegeben?



\paragraph{Vorkorrektur und Probevortrag}
Wir bieten jeweils an, die schriftliche Ausarbeitung vor der Abgabe einmal zu lesen und Korrekturen \bzw Anmerkungen dazu zu geben.
Ebenso bieten wir an, einen Probevortrag vor dem eigentlichen Seminartermin zu halten.
Hierzu ist klarzustellen, dass sich die Bewertung der Ausarbeitung sowie des Aufbau des Vortrages sich schwerpunktmäßig auf diese Vorabversion \bzw den Probevortrag bezieht, da wir letztlich nicht unsere eigenen Korrekturen bewerten wollen.
Natürlich wird aber eine entsprechende Umsetzung etwaiger Hinweise bei der Bewertung berücksichtigt, so dass eine Vorabgabe und ein Probevortrag immer vorteilhaft ist.
(\Dah eine endgültige Abgabe der Arbeit mit größeren Mängeln ist immer "`schlimmer"', als wenn diese noch nach der Vorkorrektur verbessert werden.)
